
\section{Scope and Assumptions}

The scope of this thesis is clearly delineated to focus on the active power management and economic optimization of a single residential microgrid. While the field of microgrid control is vast, encompassing aspects such as voltage stability, frequency regulation, and protection coordination, this study narrows its focus to the "tertiary control" layer specifically, the economic dispatch and energy arbitrage problem.

\subsection{System Boundaries}
The research is bounded by strict architectural constraints. The study considers a single-agent architecture, focusing on a centralized controller for a single residential dwelling equipped with a Building Energy Management System (BEMS). Interactions with neighbors or peer-to-peer (P2P) energy trading markets are considered out of scope, although the proposed DRL framework could theoretically be extended to multi-agent scenarios in future work. Furthermore, the microgrid is assumed to operate primarily in grid-connected mode. While the ability to island is a defining feature of microgrids, this thesis focuses on the interaction with the utility grid under normal operating conditions to maximize economic benefits. Consequently, islanding capability and emergency blackout resilience are not explicitly modeled in the reward function. Finally, the optimization objective is strictly financial, focusing on active power cost minimization. Reactive power support, voltage regulation, and power quality issues such as harmonic distortion are not considered. The utility grid is modeled as an infinite bus capable of absorbing or supplying any amount of active power within the physical limits of the point of common coupling (PCC).

\subsection{Modeling Assumptions}
To maintain computational tractability while ensuring realistic simulation results, several key assumptions are made. First, it is assumed that the utility's Time-of-Use (TOU) tariff structure is known 24 hours in advance. While real-time pricing (RTP) introduces price uncertainty, TOU tariffs are currently the standard for residential prosumers in many jurisdictions. Second, the battery is modeled with a constant round-trip efficiency. While internal resistance and capacity fade are non-linear processes dependent on temperature and depth of discharge (DoD), a simplified efficiency model is deemed sufficient for day-ahead economic planning. However, a degradation cost term is included in the reward function to discourage abusive cycling, simulating the long-term impact of battery health on operational decisions. Third, uncertainty is characterized specifically in the solar generation and load consumption profiles. The "Forecast" profiles provided to the agent differ from the "Actual" profiles experienced during the simulation step, mimicking real-world prediction errors. This allows for a dedicated analysis of the agent's robustness to Sim-to-Real gaps.

By adhering to these boundaries, this thesis provides a focused and rigorous evaluation of Deep Reinforcement Learning's potential to solve the specific, high-value problem of residential energy arbitrage under uncertainty.
